% Part: lambda-calculus
% Chapter: introduction
% Section: primitive-recursion

\documentclass[../../../include/open-logic-section]{subfiles}

\begin{document}

\olfileid[hi]{lam}{int}{pr}
\olsection{\usetoken{S}{lambda definable} फलन आदिम पुनरावर्तन के अधीन संवृत हैं}

आदिम पुनरावर्तन के लिए अंततः हमें कुछ काम करना पड़ेगा। हमें चरणों में आगे
बढ़ना होगा। पहले की तरह मान लें कि हमारे पास पहले से ऐसे पद $G$ और $H$ हैं
जो क्रमशः फलनों $g$ और~$h$ को !!{lambda define} करते हैं; हम ऐसा पद $H$
चाहते हैं जो निम्न प्रकार परिभाषित फलन~$f$ को !!{lambda define} करे:
\begin{align*}
f(0, \vec z) & = g(\vec z) \\
f(x+1, \vec z) & = h(z, f(x,\vec z), \vec z).
\end{align*}
अतः सामान्यतः, लैम्ब्डा पद $G'$ और~$H'$ दिए होने पर ऐसा पद~$F$ खोजना
पर्याप्त है कि
\begin{align*}
F(\num{0}, \vec z) & \equiv G(\vec z) \\
F(\overline{n+1}, \vec z) & \equiv H(\num{n}, F(\num{n}, \vec z), \vec z)
\end{align*}
प्रत्येक प्राकृतिक संख्या~$n$ के लिए; $G'$ और~$H'$ द्वारा $g$ और~$h$ को
!!{lambda define} करने का अर्थ है कि जब भी हम $\vec z$ के स्थान पर अंक
$\num{\vec m}$ रखें, तो $F(\num{n+1}, \num{\vec m})$ प्रसामान्य होकर सही
उत्तर बनेगा।

किंतु इसके लिए ऐसा पद~$F$ खोजना पर्याप्त है जो
\begin{align*}
F(\num 0) & \equiv G \\
F(\overline {n+1}) & \equiv H(\num{n},F(\num{n}))
\intertext{प्रत्येक प्राकृतिक संख्या $n$ के लिए संतुष्ट करे, जहाँ}
G & = \lambd[\vec z][G'(\vec z)] \text{ और}\\
H(u,v) & = \lambd[\vec z][H'(u,v(u,\vec z),\vec z)].
\end{align*}
दूसरे शब्दों में, लैम्ब्डा की इस युक्ति से हम अतिरिक्त प्राचलों $\vec z$
की चिंता से बच सकते हैं---वे लैम्ब्डा संकेत-लिपि में समाहित हो जाते हैं।

पद $F$ को परिभाषित करने से पहले हमें क्रमित युग्मों को सँभालने की कोई
विधि चाहिए। अगला लेम्मा यही देता है।

\begin{lem}
कोई ऐसा लैम्ब्डा पद $D$ है कि लैम्ब्डा पदों $M$ और~$N$ के प्रत्येक युग्म
के लिए $D(M,N)(\num{0}) \red M$ और $D(M,N)(\num{1}) \red N$।
\end{lem}

\begin{proof}
पहले लैम्ब्डा पद $K$ को इस प्रकार परिभाषित करें:
\[
K(y) = \lambd[x][y].
\]
दूसरे शब्दों में, $K$ पद $\lambd[y][\lambd[x][y]]$ है। दूसरे ढंग से देखें
तो प्रत्येक $M$ के लिए $K(M)$ ऐसा अचर फलन है जो किसी भी निवेश पर~$M$
लौटाता है।

अब $D(x,y,z)$ को $D(x,y,z) = z (K(y))x$ द्वारा परिभाषित करें। तब
\begin{align*}
D(M,N,\num 0) & \red \num 0 (K(N)) M \red M \text{ और}\\
D(M,N,\num 1) & \red \num 1 (K(N)) M \red K(N) M \red N,
\end{align*}
जैसा अपेक्षित था।
\end{proof}

विचार यह है कि $D(M,N)$ युग्म $\tuple{M, N}$ को निरूपित करता है; और यदि
$P$ को ऐसा युग्म निरूपित करने वाला माना जाए, तो $P(\num 0)$ और $P(\num 1)$
बाएँ और दाएँ प्रक्षेपों $(P)_0$ और $(P)_1$ को निरूपित करते हैं। हम बाद वाली
संकेत-लिपि उपयोग करेंगे।

\begin{lem}
!!{lambda definable} फलन आदिम पुनरावर्तन के अधीन संवृत हैं।
\end{lem}

\begin{proof}
हमें दिखाना है कि कोई भी पद $G$ और $H$ दिए होने पर हम ऐसा पद $F$ खोज सकते
हैं कि
\begin{align*}
F(\num{0}) & \equiv G \\
F(\num{n+1}) & \equiv H(\num{n}, F(\num{n}))
\end{align*}
प्रत्येक प्राकृतिक संख्या~$n$ के लिए। मोटे तौर पर विचार यह है कि अंकों को
पुनरावर्तक बनाकर \emph{युग्मों} के अनुक्रम
\[
\tuple{\num 0, F(\num 0)}, \tuple{\num 1, F(\num 1)}, \dots,
\]
की संगणना की जाए। ध्यान दें कि पहला युग्म केवल $\tuple{\num 0, G}$ है।
युग्म $\tuple{\num{n}, F(\num{n})}$ दिए होने पर अगला युग्म
$\tuple{\num{n+1}, F(\num{n+1})}$, ~$\tuple{\num{n+1}, H(\num{n},
F(\num{n}))}$ के तुल्य होना चाहिए। हम ऐसा लैम्ब्डा पद~$T$ बनाएँगे जो यह
एक-चरणीय संक्रमण करता है।

विवरण इस प्रकार है। $T(u)$ को परिभाषित करें:
\[
T(u) = \tuple{S((u)_0), H((u)_0,(u)_1)}.
\]
अब किसी भी संख्या~$n$ के लिए यह जाँचना सरल है कि
\[
T(\tuple{\num{n}, M}) \red \tuple{\num{n+1}, H(\num{n}, M)}.
\]
ऊपर दिए संकेत के अनुसार, $G$ और $H$ दिए होने पर~$F(u)$ को परिभाषित करें:
\[
F(u) = (u(T,\tuple{\num 0, G}))_1.
\]
दूसरे शब्दों में, निवेश~$\num{n}$ पर $F$, ~$\tuple{\num 0, G}$ पर $T$ को
$n$ बार पुनरावर्तित करता है और फिर दूसरा घटक लौटाता है। आरंभ में हमारे पास
है:
\begin{enumerate}
\item $\num{0} (T, \tuple{\num 0, G}) \equiv \tuple{\num{0}, G}$
\item $F(\num{0}) \equiv G$
\end{enumerate}
~$n$ पर आगमन द्वारा हम दिखा सकते हैं कि प्रत्येक प्राकृतिक संख्या के लिए
निम्न बातें लागू होती हैं:
\begin{enumerate}
\item $\num{n+1}(T, \tuple{\num{0}, G}) \equiv \tuple{\num{n+1}, F(\num{n+1})}$
\item $F(\num{n+1}) \equiv H(\num{n}, F(\num{n}))$
\end{enumerate}
दूसरे खंड के लिए,
\begin{align*}
F(\num{n+1}) & \red (\num{n+1}(T, \tuple{\num 0, G}))_1 \\
& \equiv (T(\num{n} (T, \tuple{\num 0, G})))_1 \\
& \equiv (T(\tuple{\num{n}, F(\num{n})}))_1 \\
& \equiv (\tuple{\num{n+1}, H(\num{n}, F(\num{n}))})_1 \\
& \equiv H(\num{n}, F(\num{n})).
\end{align*}
यहाँ अंतिम से पहले वाली पंक्ति में हमने आगमन-परिकल्पना का उपयोग किया है।
पहले खंड के लिए,
\begin{align*}
\num{n+1} (T, \tuple{\num 0, G}) &
\equiv T(\num{n} (T, \tuple{\num 0, G})) \\
& \equiv T( \tuple{\num{n}, F(\num{n})}) \\
& \equiv \tuple{\num{n+1}, H(\num{n}, F(\num{n}))} \\
& \equiv \tuple{\num{n+1}, F(\num{n+1})}.
\end{align*}
यहाँ अंतिम पंक्ति में हमने दूसरे खंड का उपयोग किया है। इस प्रकार हमने
$F(\num 0) \equiv G$ और प्रत्येक $n$ के लिए $F(\num {n+1}) \equiv H(\num
n, F(\num{n}))$ दिखा दिया, और हमें ठीक यही चाहिए था।
\end{proof}

\end{document}

